\documentclass[10pt,a4paper]{article}
\usepackage[T1]{fontenc}
\usepackage[utf8]{inputenc}
\usepackage{lmodern}
\usepackage[spanish]{babel}
\usepackage[margin=1.7cm]{geometry}
\usepackage{hyperref,fancyhdr}
\hypersetup{hidelinks}
\pagestyle{fancy}\fancyhf{}
\fancyhead[L]{Vía · EI1039 / EI1048}
\fancyhead[R]{Entrega 2 · Patrón de interfaz}
\fancyfoot[C]{\thepage}
\setlength{\headheight}{14pt}
\setlength{\parindent}{0pt}
\setlength{\parskip}{6pt}
\newcommand{\spike}[2]{\href{https://github.com/Kiliansito/ei1039-ei1048-movilidad/blob/c170fd7/entrega1/02_Spike/#1}{\texttt{#2}}}
\begin{document}
\setcounter{page}{13}
{\LARGE\bfseries MVVM en la interfaz de Vía}\par

Hemos elegido MVVM porque el planificador necesita mostrar varios estados durante la misma acción. Mientras se pide una ruta, la pantalla indica que está calculando. Después muestra el trayecto o un mensaje de error. En el spike, un \texttt{ViewModel} guarda ese estado y React lo usa para actualizar el formulario y los resultados.

\section*{Qué hace cada parte}

\textbf{Modelo.} En el servidor, \texttt{MobilityService.plan} comprueba los datos de la solicitud, consulta al proveedor de rutas y calcula las estimaciones. La API devuelve el resultado al navegador.

\textbf{Vista.} \texttt{Calculator} recoge el origen, el destino, el transporte y el criterio de ruta. \texttt{Result} muestra la alternativa seleccionada y pasa el trazado a \texttt{MapView}. La vista también guarda decisiones locales, como qué alternativa está seleccionada.

\textbf{ViewModel.} \texttt{vm.plan} solicita la planificación a la API y guarda el resultado. \texttt{vm.run} indica que hay una operación en curso y recoge sus errores. Los componentes leen el estado con \texttt{useApp}, que enlaza la suscripción del modelo con React mediante \texttt{useSyncExternalStore}.

\section*{Del formulario al mapa}

\begin{enumerate}
    \item El usuario pulsa \textbf{Calcular ruta}. El formulario comprueba que ha elegido un origen y un destino y llama a \texttt{vm.plan} con los datos.
    \item El \texttt{ViewModel} borra el resultado anterior y envía un \texttt{POST} a \texttt{/api/plan}. Para calcular las rutas que devuelve la API, el servidor usa \texttt{MobilityService.plan}.
    \item El \texttt{ViewModel} actualiza su estado y avisa a los componentes suscritos. El formulario navega a la pantalla de resultados cuando la operación termina bien. Allí, \texttt{Result} lee la ruta guardada y la muestra en el mapa. Si la petición falla, el mensaje queda en \texttt{error} y el usuario puede seguir en el formulario.
\end{enumerate}

\section*{Por qué encaja con el planificador}

El estado observable permite que el formulario, los avisos y el mapa lean el resultado de la misma operación. Un controlador de eventos en un enfoque MVC también podría coordinar la petición. En el flujo actual de React nos resulta más directo que cada vista lea el estado que necesita del \texttt{ViewModel}. Cuando ese estado cambia, React se encarga de actualizar la vista.

La separación todavía es parcial en el spike. El formulario decide cuándo navegar y también modifica directamente parte del estado del modelo con \texttt{vm.set}.

\section*{Aplicación en toda la interfaz}

Aplicaremos MVVM de forma uniforme en la interfaz de Vía, incluidas las pantallas de cuenta, lugares y planificación de rutas. Las vistas de React recogen las acciones del usuario y muestran los resultados. El \texttt{ViewModel} coordina las peticiones a la API y mantiene el estado que comparten las pantallas, mientras que las reglas de negocio se ejecutan en el servidor. Cada vista conserva el estado de sus propios controles, como la alternativa de ruta seleccionada. Usamos Observer para avisar a las vistas cuando cambia el estado del \texttt{ViewModel}.

\textbf{Código del spike.} \spike{web/main.tsx}{Vistas de React},
\spike{web/model.ts}{ViewModel},
\spike{src/http.ts}{API} y
\spike{src/services.ts}{servicio de planificación}.
\end{document}
